% ============================================================
%  SECTION 2: INSTITUTIONAL BACKGROUND
% ============================================================

\section{Institutional Background}
\label{sec:background}

\subsection{Tennessee HOPE Scholarship}

Tennessee introduced the HOPE Scholarship in 2003, funded by proceeds
from the state education lottery \parencite{THEC2025}. The program provides tuition and fee
subsidies to Tennessee residents who graduate from an eligible high
school with a minimum 3.0 GPA or score at least 21 on the ACT (or 980
on the SAT). Recipients must maintain a 2.75 cumulative GPA after the
first semester to retain eligibility
\parencite{ness2008eligibility}.

The scholarship provides a fixed dollar award covering approximately
60--80\% of in-state tuition at public institutions governed by the
Tennessee Board of Regents (TBR) and the University of Tennessee system
(\$3,000 for freshmen and sophomores, scaling to \$6,000 for juniors and
seniors at four-year institutions during the study period). Students
attending eligible private institutions receive a smaller fixed amount
(\$3,000--\$4,000 per year), which covers a substantially lower share
of private tuition. This differential generosity creates a relative price advantage
for public colleges---the key mechanism in our theoretical framework.

Table~\ref{tab:hope_distribution} reports the distribution of HOPE
recipients across sectors for the first five entering cohorts
(2004--2008). Public institutions---TBR universities, TBR community
colleges, and UT campuses---account for approximately 83--85\% of
recipients in every year. Private colleges receive 15--17\%. This
concentration of aid in the public sector motivates our identification
strategy: we treat Tennessee public colleges as the primary recipients
of the HOPE demand shock.

\subsection{Prior Research on TN HOPE}

Several studies examine Tennessee's merit aid program. The closest
antecedent is \textcite{bruce2014jackpot}, who use a regression
discontinuity at the ACT score cutoff for HOPE eligibility.
\textcite{chakrabarti2022merit} study merit aid and student mobility
more broadly, documenting how scholarship programs interact with college
selectivity. \textcite{bruce2014jackpot} find that Tennessee HOPE had no
effect on overall enrollment among Pell-eligible students near the ACT
threshold, but reallocated them from two-year to four-year public
institutions ($-2.9$ percentage points at community colleges, $+2.4$ at
four-year publics). For non-white students, overall enrollment declined
by 2.6 percentage points.

Our analysis is complementary but asks a different question. Bruce and
Carruthers estimate the \emph{individual enrollment response} at the
eligibility margin: how does a student just above the ACT cutoff choose
differently from one just below? Their RD is local to the threshold and
identifies the effect for compliers at that margin. We estimate the
\emph{institutional composition effect} across the full income
distribution: how does the parental income profile of the entire student
body at public colleges change after HOPE? Our
difference-in-differences design captures the average treatment effect
across all treated cohorts, including students well above the
eligibility threshold for whom the RD has no power. The distinction
matters because the composition of public colleges depends on the
sorting behavior of students at all income levels, not only those at
the eligibility margin. Our quintile decomposition---showing that all
five quintiles shift monotonically---documents a pattern that no
threshold-local design can detect, since students in the top quintile
are almost all well above any reasonable ACT cutoff. Their finding that
lower-income marginal students upgrade from community colleges to
four-year institutions is consistent with our result that bottom-quintile
share rises and that two-year and four-year colleges adjust through
different channels.

\subsection{Timeline and Cohort Mapping}
\label{sec:background_timeline}

The Tennessee HOPE Scholarship was enacted in 2003. In the
\textcite{chetty2020mobility} data, cohorts are indexed by birth year,
and students are assigned to institutions attended between ages 19 and
22. The analysis window spans birth cohorts 1980--1991, corresponding
approximately to college entry years 1998--2009.
Table~\ref{tab:cohort_map} maps each birth cohort to its decision year,
college entry year, event time, and HOPE exposure status.

\paragraph{Decision timing vs.\ receipt timing.}

The mechanism studied in this paper is sorting across institutional
sectors---public, private, and out-of-state---driven by the relative
price change that HOPE creates. Sorting occurs \emph{before}
matriculation, when students choose where to apply and enroll. The
economically relevant treatment date is therefore the point at which
HOPE enters the student's decision environment, not the date on which
scholarship funds are first disbursed.

We define the decision year as birth cohort $+ 17$, corresponding to
the year in which the student is typically a high-school senior forming
a college choice set. Under this mapping, the 1986 birth cohort has
decision year 2003---the same year HOPE was enacted. These students were
mid-application cycle when the program was announced; some could adjust,
but many had already committed. The 1987 birth cohort has decision year
2004 and is the first group whose entire college search occurred with
HOPE already enacted, publicized, incorporated into school counseling,
and operational in the financial aid environment.

Our preferred treatment timing is therefore $t^* = 1987$: the first
cohort whose college choice set was formed under the fully known HOPE
regime. We designate the 1986 cohort as the omitted reference period
($k = -1$), absorbing any partial treatment of this transition cohort
into the baseline. The 1985 cohort received limited retroactive benefits
through a bridge provision, but made enrollment decisions before HOPE
existed. The 1984 cohort (entering college Fall 2002) is the last
cohort with no HOPE exposure of any kind.

We test robustness to $t^* \in \{1985, 1986, 1987\}$ in
Section~\ref{sec:robustness}. The static DiD estimate is insensitive to
this choice.

\paragraph{Analysis window.}

Under the preferred timing, the sample provides five pre-treatment
cohorts (1981--1985) and five post-treatment cohorts (1987--1991), with
1986 as the omitted reference period. Each cohort's composition is
measured independently: the outcome for cohort $t$ at college $c$
reflects the parental income distribution of students born in year $t$
who attended college $c$ between ages 19 and 22. The 1980 birth cohort
is dropped from the preferred specification because Tennessee's
1980-born cohort exhibits anomalously low bottom-quintile share that
normalizes by 1982, consistent with a data artifact rather than a
substantive pre-trend. Dropping this single cohort substantially
improves sensitivity to violations of parallel trends
(Section~\ref{sec:robustness}).

\begin{table}[H]
\centering
\begin{threeparttable}
\caption{Cohort Mapping: Birth Year, Decision Year, and HOPE Exposure}
\label{tab:cohort_map}
\begin{tabular}{cccclc}
\toprule
Birth Cohort & Decision Year & College Entry & Event Time ($k$) & HOPE Status \\
\midrule
1980 & 1997 & Fall 1998 & --- & No HOPE (dropped$^\dagger$) \\
1981 & 1998 & Fall 1999 & $-6$ & No HOPE \\
1982 & 1999 & Fall 2000 & $-5$ & No HOPE \\
1983 & 2000 & Fall 2001 & $-4$ & No HOPE \\
1984 & 2001 & Fall 2002 & $-3$ & No HOPE \\
1985 & 2002 & Fall 2003 & $-2$ & Bridge provision (partial)$^\ddagger$ \\
1986 & \textbf{2003} & Fall 2004 & $-1$ & Transition (ref.\ period) \\
\midrule
1987 & \textbf{2004} & Fall 2005 & $0$ & Full HOPE ($t^*$) \\
1988 & 2005 & Fall 2006 & $+1$ & Full HOPE \\
1989 & 2006 & Fall 2007 & $+2$ & Full HOPE \\
1990 & 2007 & Fall 2008 & $+3$ & Full HOPE \\
1991 & 2008 & Fall 2009 & $+4$ & Full HOPE \\
\bottomrule
\end{tabular}
\begin{tablenotes}[flushleft]
\footnotesize
\item \textit{Notes:} Decision year $=$ birth cohort $+ 17$ (high-school
  senior year). College entry is approximate (birth cohort $+$ 18).
  Students are assigned to institutions attended between ages 19 and 22
  in the \textcite{chetty2020mobility} data. HOPE was enacted in 2003
  (bolded). $^\dagger$The 1980 cohort is dropped from the preferred
  specification; see text. $^\ddagger$Bridge provision extended limited
  retroactive benefits to Fall 2003 entrants who met eligibility
  criteria; these students made enrollment decisions before HOPE existed.
  Event times shown relative to $t^* = 1987$.
\end{tablenotes}
\end{threeparttable}
\end{table}

\paragraph{Confounding policies.}

Several other policy changes coincided with HOPE's introduction, and we
consider whether any could produce the quintile gradient we document.

Tennessee's education lottery launched in 2003 and generated revenue
that funded both HOPE and Pre-K programs. The federal No Child Left
Behind Act (2002) altered K--12 accountability requirements nationwide.
Both policies affect treated and control states symmetrically and are
absorbed by cohort fixed effects.

Federal Pell Grant generosity increased during our analysis window:
the 2007 College Cost Reduction Act raised maximum awards, and the
2009 American Recovery and Reinvestment Act added temporary increases.
These expansions apply nationally, so any composition effects are
absorbed by cohort fixed effects in our DiD design. To the extent
that Pell expansions disproportionately benefit low-income students,
they would produce a level shift in bottom-quintile enrollment at all
public colleges, not the Tennessee-specific treatment effect our
interaction term identifies.

Tennessee's Board of Regents expanded community college access during
this period, including the Tennessee Technology Centers and dual
enrollment initiatives. These reforms could affect overall enrollment
levels but would not generate a monotonically ordered gradient across
all five income quintiles---a pattern that requires differential price
sensitivity by parental income, as our model predicts.

Tuition changes at neighboring states' public universities could
affect cross-state sorting. However, such changes would
disproportionately influence border-county colleges rather than
producing the statewide pattern we observe.

While we cannot rule out every potential confounder individually, the
quintile gradient provides a useful discriminating pattern: a generic
enrollment shock would affect quintile shares unevenly but would be
unlikely to produce five coefficients that are monotonically ordered
across all quintiles with no effect at private institutions.

\subsection{Georgia HOPE as a Comparison Case}
\label{sec:background_georgia}

Georgia's HOPE Scholarship (1993) is the canonical case in the merit aid
literature \parencite{dynarski2000hope, song2024hope} and our primary
theoretical comparison. Three design features
distinguish it from Tennessee's program in ways that matter for our
sorting framework.

First, from 1993 to 2000, Georgia HOPE awards were reduced
dollar-for-dollar by Pell Grant amounts. Students from the
lowest-income families received zero net benefit from HOPE for the
program's first seven years. \textcite{dynarski2000hope} documents the
consequence: Black students, who were disproportionately Pell-eligible,
showed no attendance gain, while white students gained 12.3 percentage
points. The Pell offset effectively eliminated the access margin for
the bottom of the income distribution.

Second, Georgia imposed an income cap (\$66,000, later \$100,000,
eliminated in 1995) during the program's early years. While the cap
was removed relatively quickly, its presence during the launch shaped
initial enrollment patterns.

Third, Georgia HOPE provided comparable awards to private colleges:
initially full tuition, later a fixed \$3,000 that covered a
meaningful share of private tuition at Georgia's price levels. This
smaller public-private subsidy differential muted the sector-switching
mechanism that drives compositional sorting in our model. \textcite{cornwell2006enrollment} confirm that four-year private enrollment in
Georgia \emph{rose} by 13.2\% after HOPE, and two-thirds of the
four-year enrollment gain reflected reduced out-of-state migration
rather than new college-goers.

We return to the GA-TN comparison in Section~\ref{sec:theory_georgia},
where we show that these design differences generate opposite
predictions from the same theoretical framework.

\subsection{Florida Bright Futures: A Tuition-Scaling Contrast}
\label{sec:background_florida}

Florida's Bright Futures Scholarship (1997) provides a structural
contrast to Tennessee's fixed-dollar design. Bright Futures operates
on a multi-tier system. The top tier (Florida Academic Scholars)
requires a 3.5 weighted GPA and 1270~SAT, covering 100\% of tuition.
The broad-coverage tier---Florida Medallion Scholars---requires only
a 3.0 weighted GPA and 970~SAT (old scale, approximately the 50th
percentile nationally) and covers 75\% of tuition at public
institutions. During our study period (1998--2009), approximately
30--35\% of Florida high school graduates received some tier of
Bright Futures.

Two features distinguish Bright Futures from Tennessee HOPE for the
income-shock analysis. First, the award is \emph{percentage-based}:
it covers a fixed share of tuition rather than a fixed dollar amount.
When tuition rises, the award rises proportionally. This means the
gap between the scholarship and total tuition cost does not widen
over time---the theoretical distinction formalized in
Proposition~\ref{prop:compensation}. Second, Bright Futures does not
require FAFSA filing for eligibility, removing a paperwork barrier
that may deter low-income families during periods of financial stress.

Because Bright Futures was already active throughout the Chetty panel
(the program predates the 1998 college entry of the earliest cohort),
we cannot estimate an adoption DiD for Florida. The shock-interaction
design (Section~\ref{sec:empirical_shock}) circumvents this limitation
by comparing how Florida colleges respond to income shocks relative
to non-merit states, without requiring a pre-program period.
